%% kapa: exam=1 round=37 year=2026 period=1 form=A subject=civil_law subject_ko=민법 seq=16 number=16
%% source: Q-Net 감정평가사 기출문제 (cst00309) artlSeq=5258924 file=raw/1008/5258924/2026년 제37회 감정평가사 1차 1교시.pdf
%% conversion: kordoc 4.24.51 -> tools/gen1.py (auto); review=proofread-sonnet; flags=none
\begin{question}{16}
법률행위의 무효에 관한 설명으로 옳지 않은 것은? (다툼이 있으면 판례에 따름)
\begin{choices}
  \item 무효인 법률행위의 추인은 묵시적인 방법으로도 할 수 있다.
  \item 의사무능력으로 인한 무효행위의 추인은 그 무효원인이 소멸한 후에 하여야 효력이 있다.
  \item 무효인 법률행위에 따른 법률효과를 침해하는 것처럼 보이는 위법행위가 있는 경우, 그 법률효과의 침해에 따른 손해배상을 청구할 수 있다.
  \item 무권리자의 처분이 계약으로 이루어진 경우에 권리자가 이를 추인하면 원칙적으로 계약의 효과가 계약을 체결하였을 때에 소급하여 권리자에게 귀속한다.
  \item 토지거래허가구역 내 토지에 관하여 허가를 잠탈하는 내용으로 체결된 매매계약은 그 후 허가구역 지정이 해제되더라도 유효로 되지 않는다.
\end{choices}
\end{question}
\answer{3}
